% /root/monodromy, 2026-11-19. Complete independent root, RD and OC audits passed.
\subsection{A point-incidence criterion for excluding all other kernel collisions}
\label{subsec:ss-uniform-collision-exclusion}
\status{Proved. Three independent complete proof audits and exact
program audits passed. This removes the other-collision gate for a
prescribed mixed contact under every finite odd kernel with injective
coordinate projections. Uniform tangent-triple bridges are established
in Proposition~\ref{prop:ss-uniform-tangent-triple}; the full image
in every odd degree follows in Theorem~\ref{thm:ss-all-odd-full-groups}.}

Work over an algebraically closed field $k$ of characteristic two.
Let $E:y^2+y=x^3$, $X=E^2$ and
$L=\mathcal O_E(3O)\boxtimes\mathcal O_E(3O)$.
Let $K\subset X(k)$ be a finite subgroup of odd order such that both
projections $K\to E(k)$ are injective. Fix $0\ne g\in K$, written
$g=((a,b),(c,d))$. In particular both coordinates are nonzero.
Use the homogeneous contact slice $\Lambda=\mathbf P^4$ with coordinates
$[F:T:U:v:h]$ from
\eqref{eq:ss-contact-linear-system}--\eqref{eq:ss-contact-cubic-homogeneous}.
We also use that slice when $g$ has order three. Write
$\mathcal Q=(q=0)\subset\Lambda$ and $P=(O,O)$, $Q=g$.

\begin{lemma}[The contact cubic and its evaluation map]
\label{lem:ss-contact-cubic-evaluation}
The cubic $\mathcal Q$ is smooth. It contains no projective plane.
Its intersection with every projective hyperplane is an integral
cubic surface. For each $0\ne j\in K$, the locus where the two value
functionals at $R,R+j$ on $\Lambda$ have rank at most one is finite.
Their common rank is zero only when $j=\pm g$, at the corresponding
ordered pair of base points $P,Q$. No three distinct points in a
$K$-orbit are all base points of $\Lambda$.
\end{lemma}
\begin{proof}
Set $A=a^4+a$ and $B=c^4+c$. In the ordered variables $(F,T,U,v,h)$,
the cubic $q=z^{[2]t}A_qz$ has coefficient matrix
\[
 A_q=\begin{pmatrix}
 0&1&0&0&A\\
 1&0&0&0&0\\
 0&0&1&0&0\\
 0&0&0&0&1\\
 0&0&B&1&0
 \end{pmatrix},\qquad \det A_q=1.
\]
Its gradient is $A_q^t z^{[2]}$, so it has no projective singular
point. If a vector subspace $W$ is contained in its affine cone,
expanding $q(u+t w)$ gives
$u^{[2]t}A_qw=0$ for all $u,w\in W$.
Consequently $2\dim W\le5$. This excludes projective planes.
A reducible cubic in a projective hyperplane $\mathbf P^3$ would
have a linear factor and hence contain a projective plane. Therefore
every such section is integral.

For the value-map calculation put
\[
 X_1=x+a,\quad Y_1=y+b,\quad Z_1=z+c,\quad
 \ell=w+c^2z+d^2,\quad P_1=Y_1+a^2X_1.
\]
The section can be rewritten as
\begin{equation}\label{eq:ss-contact-five-functions}
       s=(T+hX_1)\ell+FP_1+vX_1+UX_1Z_1.
\end{equation}
Equivalently the five functions defining the rational value map are
$X_1,Y_1,\ell,X_1z,X_1\ell$.
Where both factors are affine and $X_1\ell\ne0$, their ratios
recover $z,\ell,X_1,Y_1$, and therefore the point of $X$.
The images of points outside this open lie outside its image: in
these coordinates their first or fifth homogeneous coordinate is zero.

It remains to examine the exceptional fibres. The three vertical
fibres through $O,g_1,-g_1$ have respectively constant image $F$,
constant image $T$, and nonconstant image on the line spanned by
$F,T$. The horizontal fibre through $O$ maps nonconstantly to the
line spanned by $T,h$. The zero divisor of $\ell$ has support
$g_2,-2g_2$; its two horizontal fibres have nonconstant images on
the lines spanned respectively by $F,v$ and
$F,v+\bigl(x(-2g_2)+c\bigr)U$.
For a nonflex $g_2$ these are distinct lines. For a flex the two
horizontal fibres coincide, and there is just the line spanned by
$F,v$. These assertions follow directly from
\eqref{eq:ss-contact-five-functions}, including the infinity
charts, where the section is divided by the local trivializing
monomial of $\mathcal O(3O)$.

A nonzero element of $K$ translates both elliptic coordinates
nontrivially. Thus an exceptional vertical or horizontal fibre cannot
be translated to itself. Distinct fibres of the same type above
either have distinct constant images, or a constant image and a
nonconstant image, or nonconstant images on distinct lines.
In every case equality of the two values occurs at only finitely
many points. This proves the rank-one assertion; any positive
dimensional equality locus would have had to be such an exceptional
fibre.

The base locus is $P,Q,(g_1,-2g_2)$, with the last two points
coinciding when $3g=0$. Injectivity of the first projection of $K$
shows that the third point has no $K$-difference with either of the
first two unless it is $Q$. This proves both remaining assertions.
\end{proof}

\begin{lemma}[The moving two-point covers are generically separable]
\label{lem:ss-contact-pair-covers}
For a geometric generic section of $\mathcal Q$, the intersections
of its divisor with all its nontrivial $K$-translates are finite.
For $j\ne\pm g$, every intersection with its $j$-translate is
transverse. For $j=g$, there is the prescribed intersection at $P$;
its length is two if $3g\ne0$ and three if $3g=0$. All other
intersections with the $g$-translate are transverse. The assertion
for $-g$ is the translated version of this one.
\end{lemma}
\begin{proof}
Consider the proper incidence
\[
 I_j=\{(R,s)\in X\times\mathcal Q:s(R)=s(R+j)=0\}.
\]
On its generic rank-two value locus over $X$, the coefficient fibre
is a plane cubic, never the whole plane by the preceding lemma.
We first prove that this generic plane cubic is smooth.

Suppose $j\ne\pm g$. At $R=P$ the first value is zero and the
second is a nonzero functional $l_j$. Since $j\ne\pm g$,
$X_j=x(j_1)+a\ne0$. Put
\[
 r=\ell(j)/X_j,\qquad p=P_1(j)/X_j,\qquad Z_j=x(j_2)+c.
\]
The equation $l_j=0$ is
$v=rT+\ell(j)h+pF+Z_jU$.
In the two $Y=1$ coordinates at $P$, approach $P$ along
$(t_1,t_2)=(\kappa t^3,t)$. The first value, divided by $t^3$,
tends to $F+\kappa h$. Hence the limiting coefficient plane is
defined by $l_j=F+\kappa h=0$. Substitution into $q$ shows that
its Frobenius-form coefficient matrix, in variables $(T,U,h)$,
has determinant
\begin{equation}\label{eq:ss-pair-plane-determinant}
                        (\kappa+r^2)(\kappa^2+r).
\end{equation}
Choose $\kappa$ outside the finite zero set of this expression.
The limiting plane cubic is smooth, so the geometric generic
plane cubic is smooth and integral.

For $j=g$, instead approach $P$ along $(t,0)$. The two values,
divided by $t$, tend to $h,v$. The limiting plane cubic is
\[
                           U^3+F^2T+FT^2=0,
\]
again smooth. This proves the generic integrality in this case too.

Consequently each $I_j$ has exactly one three-dimensional component
dominating $X$, which is generically reduced. All rank-one strata
are finite over $X$ and have two-dimensional coefficient fibres.
For $j\ne\pm g$ there is no rank-zero stratum, hence no other
component can dominate $\mathcal Q$. For $j=g$ the only additional
three-dimensional support is the constant section $R=P$.

For $j\ne\pm g$ we exhibit an etale point of $I_j\to\mathcal Q$.
Choose $s$ in the integral cubic surface $\mathcal Q\cap(l_j=0)$
such that $h\ne0$ and $\partial_2s(j)\ne0$.
Both are genuine nonzero linear conditions on $\ker l_j$:
$l_j$ has nonzero $v$ coefficient $X_j$, whereas
\[
 \partial_2s(j)=(T+hX_j)\bigl(x(j_2)^2+c^2\bigr)+UX_j
\]
has no $v$ coefficient and has nonzero $U$ coefficient.
The Jacobian of the two value equations with respect to $R$ at
$R=P$ has determinant $h\partial_2s(j)\ne0$.
This point therefore belongs to the unique main component, whose
projection is dominant and generically etale. After removing a
proper closed subset of $\mathcal Q$, its entire finite cover is
etale; all lower-dimensional components have also disappeared.

For $j=g$ and $3g\ne0$, the same argument at $R=Q$ applies to
the other value at $2g$. Here $2g\ne\pm g$ and
$x(2g_2)+c\ne0$. The integral surface
$\mathcal Q\cap(s(2g)=0)$ is not contained in $v=0$ or in
$\partial_2s(2g)=0$. The Jacobian determinant is
$v\partial_2s(2g)$, so the main component is again generically
etale. The constant intersection at $P$ has length two by
Lemma~\ref{lem:ss-uniform-local-contact}; its leading coefficient
is $T(c^4+c)$.

If $3g=0$, the constant intersection has length exactly three.
Indeed the branch through $P$ has first coordinate
$(F/h)t^3+O(t^4)$, and translation on the second elliptic factor
gives $\ell(g_2+T(t))=t^3+O(t^4)$. The other value on that branch
therefore has leading coefficient
\begin{equation}\label{eq:ss-flex-contact-length}
                             T+vF/h.
\end{equation}
It is generically nonzero on the integral cubic $\mathcal Q$.
Every component of $I_g$ has dimension at most three, so its
generic fibre over $\mathcal Q$ is finite: a positive-dimensional
fibre over the generic point would require a component of dimension
at least four. Intersection theory then gives its total length
$L^2=18$. Subtracting the length-three constant component leaves
degree $15$ on the unique integral main component. Its inseparable
degree is a power of two dividing $15$, so it is separable.
After shrinking $\mathcal Q$ this moving cover is etale.
The same dimension argument and intersection number apply in the
preceding cases. This proves every assertion.
\end{proof}

\begin{lemma}[There are generically no three kernel translates]
\label{lem:ss-contact-no-triples}
A geometric generic section of $\mathcal Q$ contains no three
distinct points of a $K$-orbit.
\end{lemma}
\begin{proof}
For distinct nonzero $j,k\in K$, consider the three-value incidence
over $X$. We first show that its generic three value functionals
are independent and that $q$ does not vanish identically on their
common vector kernel of dimension two.

One may reorder and translate the unordered triple. Unless it is
a coset of $\{0,g,-g\}$ with $3g=0$, choose its initial point so
that neither $j$ nor $k$ equals $g$. This is possible because a
three-element set in which every point has its $g$-successor in the
set is a three-element orbit of translation by $g$.
At $R=P$, the other two value rows $l_j,l_k$ are independent modulo
$F,h$. For $j\ne\pm g$ their rows in $(T,v,U)$ are
\[
                        (\ell(j),X_j,X_jZ_j).
\]
Different second-factor $x$-coordinates give different normalized
last entries. Equal second-factor $x$-coordinates imply $k=-j$,
by injectivity of the second projection of $K$; then $X_k=X_j$
and $\ell(k)=\ell(j)+1$, so the rows are still independent.
The exceptional allowed index $j=-g$ has row $(1,0,0)$, also
independent from every other permitted row. Thus
$l_j,l_k,F,h$ have rank four.

Approach $P$ along $(\kappa t^3,t)$ as before. The limiting three
rows are $F+\kappa h,l_j,l_k$, of rank three. If the cubic vanished
on the generic common kernel, it would vanish on every limiting
kernel
\[
                    \ker(l_j,l_k,F+\kappa h).
\]
As $\kappa$ varies, these two-dimensional vector spaces have dense
union in the three-dimensional space $\ker(l_j,l_k)$, because
$F,h$ are independent on that space. This would put a projective
plane inside $\mathcal Q$, contrary to
Lemma~\ref{lem:ss-contact-cubic-evaluation}.

For the exceptional triple take $(j,k)=(g,-g)$ and $3g=0$.
Then $a^4+a=c^4+c=0$. Along the same path, after division of the
first two rows by $t^3$, the limiting rows are
\[
                      F+\kappa h,\quad T+\kappa v,\quad F+T.
\]
The coefficient one in the middle row is the actual elliptic
translation coefficient
$\ell(g_2+T(t))=t^3+O(t^4)$.
For $\kappa\ne0$ the three rows have rank three, and their kernel
has $F=T=\kappa h$, $v=h$. The cubic restricts to $U^3$, which
is not zero. This proves the required generic assertion in all cases.

The incidence dimension now proves the lemma. On the generic
rank-three stratum the cubic cuts the coefficient line in finitely
many points. Where it contains that line, the base has dimension at
most one, so this stratum still has total dimension at most two.
The rank-two locus has dimension at most one, and its coefficient
plane is never contained in $\mathcal Q$, again giving total
dimension at most two. The rank-one locus is finite by the pair
calculation, with coefficient fibres of dimension two. There is
no rank-zero triple. Thus the full proper three-value incidence
has dimension at most two. Its image cannot fill the
three-dimensional cubic $\mathcal Q$. There are only finitely many
ordered pairs $j,k$, so they can all be excluded simultaneously.
\end{proof}

\begin{proposition}[A mixed contact with no other collisions for every graph kernel]
\label{prop:ss-uniform-graph-mixed-contact}
For every finite odd subgroup $K\subset E^2(k)$ with injective
coordinate projections, and every $0\ne g\in K$, there is a smooth
source $C\in|L|$ on the source cubic $Q_H=0$ and outside $\det S=0$
whose image in $X/K$ has exactly one simple self-tacnode, with
preimages differing by $g$, and otherwise only ordinary nodes.
The extra source Cartier-kernel differential has unequal values on
the two identified tangent vectors. In particular the local
mixed-contact criterion applies to this norm image.
\end{proposition}
\begin{proof}
We first record the irreducibility needed to combine separate open
conditions. Let $\mathcal J_g$ parametrize triples $(R,\tau,s)$,
where $\tau$ is a common translated tangent direction at $R,R+g$,
and the section vanishes at both points and has derivative zero
along $\tau$ at both points. The base $(R,\tau)$ has dimension
three. Off finitely many $R$, the six first jets at $R,R+g$ are
independent: the tensor tangent-space calculation reduces a failure
to the crossed conditions
$3R_i+g_i=0$ or $3R_i+2g_i=0$ on both elliptic factors.
Thus on this open base the four contact conditions are independent,
and its coefficient spaces form a rank-five vector bundle.
The full source cubic $Q_H$ is a nondegenerate Frobenius form in
nine variables. It contains no vector space of dimension five,
by the same polarization argument as above. Consequently every
coefficient fibre is a genuine cubic threefold.
At $(R,\tau)=(P,\text{second-factor tangent})$ the fibre is the
smooth cubic $\mathcal Q$. The six jets here are indeed independent:
the origin rows select coefficients $S_{22},S_{12},S_{21}$, and
the three translated rows have a unit triangular minor on
$S_{00},S_{10},S_{01}$.
The relative cubic is nonzero in every fibre, so it is a flat
relative hypersurface in the projective bundle. Its constant fibre
dimension three excludes vertical components. The smooth integral
fibre just exhibited therefore gives a geometrically integral
generic fibre, and the total family on the good base is integral.
There is a unique six-dimensional component of
$\mathcal J_g\cap(Q_H=0)$.

There is no further component of that dimension over the finitely
many exceptional $R$. Indeed the two values and the derivative in
any nonzero tangent direction at the first point have rank at
least three. To see this, the Segre point representing $R+g$
does not lie in the embedded tangent space at $R$: both its factor
coordinates differ from those of $R$, and its image in the tensor
product of the two one-dimensional-point quotients is nonzero.
Thus the contact rank is at least three. The coefficient vector
space then has dimension at most six; intersecting with $Q_H$
reduces its projective dimension to at most four. Together with
the one-dimensional choice of tangent, this gives dimension at
most five. We work on the unique main component from now on.

The preceding two lemmas give, on its particular vertical slice,
a section with no triple and with no unwanted nontransverse
identification. All pair intersections are finite. These conditions
give a nonempty open on the main contact component. The only
openness point needing care is $3g=0$, when the prescribed
intersection in the slice has length three, not two. Near a chosen
point the proper intersection family is finite flat of degree
$18$: the two Cartier equations form a regular sequence in every
fibre of the smooth relative surface, so their zero-dimensional
complete intersection is flat over the base. Its length-three cluster at the marked point
contains the prescribed intersection of length at least two.
A further tangency in that cluster would contribute another
length at least two, requiring total length at least four.
It is therefore impossible. The finite-cluster decomposition is
made etale locally on the coefficient base; the resulting open
condition descends to the base. The other clusters are etale and
remain so after shrinking. This proves the stated nonempty open also in
the flex case. The triple-avoidance conditions are open directly
by their proper incidences. For $j\ne\pm g$, transversality is
likewise an open condition by the proper two-value incidence.

There is also a nonempty open of this same main component on
which the source is globally smooth, the prescribed contact is
simple, $\det S\ne0$, and the two Cartier values are unequal.
If $3g\ne0$, these assertions follow from
Lemmas~\ref{lem:ss-uniform-local-contact} and
\ref{lem:ss-uniform-source-smoothness}; injectivity of the two
projections makes both coordinates have the same non-three order.
If $3g=0$, product automorphisms of $E$ preserving $O$ take the
already certified contact of Lemma~\ref{lem:ss-global-contact03}
to the prescribed nonzero pair of three-torsion coordinates.
These automorphisms preserve the linear system and the two source
Hasse factors. All asserted conditions are open on the main
incidence. At the actual two points of contact03 the six jets have
rank six, as independently certified by
\path{agents/monodromy/code/verify_ss_collision_main_incidence.py},
with the six rows and their row reduction retained in its output.
Thus this witness, as well as the entire vertical slice at $P$,
belongs to the unique main incidence.

The two nonempty opens intersect by irreducibility. For a section
in their intersection, all translates meet in finitely many
points, so the map of its smooth source to $X/K$ is birational
onto its image. The absence of triples and the transverse
intersection conclusions give precisely ordinary nodes away from
the one simple tacnode. The unequal-value assertion is exactly
the hypothesis of the already proved local mixed-contact Cartier
criterion. This completes the proof.
\end{proof}

\begin{remark}[Scope and exact algebra certificate]
The only group restriction above is injectivity of both coordinate
projections. In particular it includes balanced odd quaternionic
isogenies with kernel $\{(-bP,P):P\in E[d]\}$ and
$b^\dagger b=d-1$. The exact program
\path{agents/monodromy/code/ss_uniform_collision_identities.py}
checks the homogeneous coefficient matrix, its determinant, the
plane-cubic determinant, the two flex limiting equations and the
order-three elliptic translation coefficient. It does not replace
the incidence, flatness or openness arguments. The subsequent uniform
tangent-triple theorem merges every possible triple-collision component
of these norm families; its odd-degree full-group consequence is
Theorem~\ref{thm:ss-all-odd-full-groups}. The full images in all even
degrees remain a separate question.
\end{remark}
